\documentclass[a4paper,11pt]{article}

\usepackage[a4paper,margin=2cm]{geometry}
\usepackage[T1]{fontenc}
\usepackage[utf8]{inputenc}
\usepackage[ngerman,english]{babel}
\usepackage{lmodern}
\usepackage{microtype}
\usepackage{xcolor}
\usepackage{parskip}
\usepackage{enumitem}
\usepackage[most]{tcolorbox}
\usepackage{titlesec}
\usepackage[hidelinks]{hyperref}
\usepackage{array}
\usepackage{longtable}
\usepackage[table]{xcolor}

\definecolor{HeaderBlueDark}{RGB}{70,110,170}
\definecolor{RowBlueLight}{RGB}{235,243,255}
\definecolor{RowBlueDark}{RGB}{220,233,250}
\definecolor{SoftGray}{RGB}{90,90,90}

\setlength{\parindent}{0pt}
\setlist[itemize]{leftmargin=1.4em,itemsep=0.35em,topsep=0.35em}
\linespread{1.06}
\renewcommand{\arraystretch}{1.35}
\setlength{\tabcolsep}{6pt}

\titleformat{\section}[block]
{\Large\bfseries\color{HeaderBlueDark}}
{}{0pt}{}
\titlespacing{\section}{0pt}{1.3em}{0.8em}

\titleformat{\subsection}[block]
{\large\bfseries\color{HeaderBlueDark}}
{}{0pt}{}
\titlespacing{\subsection}{0pt}{1.0em}{0.6em}

\newtcolorbox{exbox}{enhanced,breakable,colback=RowBlueLight,colframe=RowBlueDark,boxrule=0.4pt,arc=1.5mm,left=1.2mm,right=1.2mm,top=1mm,bottom=1mm,title={Beispiele \textnormal{(Examples)}},fonttitle=\bfseries,colbacktitle=RowBlueDark,coltitle=black}

\NewDocumentEnvironment{grammarentry}{mm+b}{%
    \begin{tcolorbox}[enhanced,breakable,colback=white,colframe=HeaderBlueDark,boxrule=0.6pt,arc=1.5mm,left=1.5mm,right=1.5mm,top=1mm,bottom=1mm,colbacktitle=HeaderBlueDark,coltitle=white,fonttitle=\bfseries,title={#1\hfill\normalfont\footnotesize #2}]
    #3
    \end{tcolorbox}
}{}

\newcommand{\GermanText}[1]{\textbf{Deutsch: }#1\par}
\newcommand{\EnglishText}[1]{{\small\color{SoftGray}\textbf{English: }#1\par}}
\newcommand{\ExampleItem}[2]{\item \textbf{#1}\\{\small\color{SoftGray}#2}}

\title{Verlaufsform der Zukunft\\\large\textnormal{Future Continuous in German}}
\date{}

\begin{document}
    \maketitle
    \tableofcontents
    \newpage


\section{Grundidee \textnormal{(Basic idea)}}
\begin{grammarentry}{Verlauf statt eigener Tempusform}{No separate German continuous tense}
\GermanText{Englisches will be + -ing wird meist durch Präsens mit Zukunftsangabe oder durch Futur I mit gerade / dabei sein ausgedrückt. Ein eigenes Futur Continuous gibt es nicht.}
\EnglishText{English will be + -ing is commonly expressed through present tense with a future time expression or through Futur I with gerade / dabei sein. German has no separate future progressive tense.}
\end{grammarentry}

\section{Future Continuous \textnormal{(English future continuous)}}
\begin{grammarentry}{Zukunftsbezug plus Verlauf}{No separate future progressive tense}
\GermanText{Auch das englische will be + -ing besitzt im Deutschen kein eigenes Tempus. Häufig genügt Präsens mit einer eindeutigen zukünftigen Zeitangabe; alternativ ist Futur I mit gerade oder einer Verlaufsumschreibung möglich. Das Futur I allein garantiert keine Verlaufsbedeutung.}
\EnglishText{German likewise has no separate tense corresponding to English will be + -ing. Present tense with a clear future time expression is often sufficient; future tense with gerade or an explicit progressive paraphrase is also possible. Future tense alone does not guarantee an ongoing interpretation.}
\end{grammarentry}
\begin{exbox}
\begin{itemize}
\ExampleItem{Morgen um zehn arbeite ich.}{Tomorrow at ten I will be working / I will work (context decides).}
\ExampleItem{Morgen um zehn werde ich gerade arbeiten.}{Tomorrow at ten I will be working.}
\ExampleItem{Morgen um zehn werde ich gerade dabei sein, einen Bericht zu schreiben.}{Tomorrow at ten I will be in the middle of writing a report.}
\ExampleItem{Morgen um zehn bin ich noch am Arbeiten.}{Tomorrow at ten I will still be working (am-progressive, depending on register).}
\end{itemize}
\end{exbox}

\section{Zukünftiger Verlauf im Passiv \textnormal{(Future ongoing passive)}}
\begin{grammarentry}{werden + Partizip II}{Passive, not automatically future}
\GermanText{Mit einer zukünftigen Zeitangabe kann das Präsens des Vorgangspassivs einen Vorgang ausdrücken, der dann gerade läuft. Werden ist hier das Passivhilfsverb.}
\EnglishText{With a future time expression, the present event passive can describe a process that will be ongoing then. Here werden is a passive auxiliary.}
\end{grammarentry}
\begin{exbox}
\begin{itemize}
\ExampleItem{Morgen um zehn wird der Bericht gerade überprüft.}{Tomorrow at ten the report will be under review / will be being checked.}
\end{itemize}
\end{exbox}
\section{Zusammenfassung \textnormal{(Summary)}}
\begin{grammarentry}{Merksätze}{Key rules}
\GermanText{Morgen um zehn arbeite ich und Morgen um zehn werde ich gerade arbeiten können beide zukünftigen Verlauf beschreiben. Gerade betont den laufenden Vorgang; die Zeitangabe verankert ihn in der Zukunft.}
\EnglishText{Morgen um zehn arbeite ich and Morgen um zehn werde ich gerade arbeiten can both describe future ongoing activity. Gerade emphasizes the process; the time expression establishes the future reference point.}
\end{grammarentry}

\section{Quelle \textnormal{(Source)}}
        \GermanText{Institut für Deutsche Sprache (IDS), grammis: Darf man Ich bin am Schreiben schreiben?}
        \EnglishText{Institute for the German Language (IDS), grammis: discussion of progressive constructions and their formation in different tenses.}
        \url{https://grammis.ids-mannheim.de/fragen/4551}
\end{document}
